% Originally: content/week-12.tex, \section{Integration of differential forms on submanifolds} (Corsin Week 12)

\section{Integration of differential forms on submanifolds}
\label{sec:integration_of_forms}

% Source: Corsin Nick/Class Notes/Week 12.pdf, p. 6
If $\mathcal{M}$ is an $m$-dimensional submanifold of $\mathbb{R}^m$ (i.e.\ simply an open set)
and $\omega \in \Omega^m(\mathcal{M})$ is an $m$-form, then
\[\int_{\mathcal{M}} \omega := \int_{\mathcal{M}} \omega_x(e_1,\ldots,e_m)\,dx;\]
in particular, for $f \in C^\infty(\mathcal{M})$,
\[\int_{\mathcal{M}} f\,dx^1\wedge\cdots\wedge dx^m = \int_{\mathcal{M}} f(x)\,dx,\]
recovering the ordinary (Riemann) integral of $f$.

\begin{definition}[Integral of a form over a parametrized submanifold]
\label{def:integral_of_form_over_submanifold}
If $L^k \subseteq \mathbb{R}^n$ is a $k$-dimensional submanifold, $n>k$, and $\phi : U \to L$
is a parametrization with $U \subseteq \mathbb{R}^k$ open, then the integral of $\omega \in
\Omega^k(L)$ is defined as
\[\int_L \omega := \int_U \phi^*\omega = \int_U \omega_{\phi(x)}\!\left(\frac{\partial\phi}
  {\partial x^1},\ldots,\frac{\partial\phi}{\partial x^k}\right) dx.\]
\end{definition}

In particular, for $f \in C^\infty(L)$ and $1\leq i_1<\cdots<i_k\leq n$,
\[\int_L f\,dx^{i_1}\wedge\cdots\wedge dx^{i_k} = \int_U f(\phi(x))\det\!\left(\frac{\partial
  \phi^{i_s}}{\partial x^t}\right)_{1\leq s,t\leq k} dx,\]
because, by property \textbf{(d)} of \cref{def:wedge_product},
\[dx^{i_1}\wedge\cdots\wedge dx^{i_k}(\partial_1\phi,\ldots,\partial_k\phi) = \det\!\left(
  \begin{smallmatrix} dx^{i_1}(\partial_1\phi) & \cdots & dx^{i_1}(\partial_k\phi) \\
  \vdots & \ddots & \vdots \\ dx^{i_k}(\partial_1\phi) & \cdots & dx^{i_k}(\partial_k\phi)
  \end{smallmatrix}\right) = \det\!\left(\frac{\partial\phi^{i_s}}{\partial x^t}\right)_{1\leq
  s,t\leq k}.\]

% Source: Linus Lüchinger/Slides/Slides-05-11.pdf, slides 7--8
% Generator: Claude Opus 5 (High)
\begin{aiexample}[Evaluating a $2$-form on a concrete pair of vectors]
\label{ex:evaluating_a_two_form}
Everything above is stated symbolically, and the mechanics are easy to lose track of. Here is
the whole computation on numbers. It is worth doing once, because every integral in this chapter
is an integral of exactly this evaluation.

The recipe implicit in the displays above is: \textbf{a wedge $dx^{i_1}\wedge\cdots\wedge
dx^{i_k}$ evaluated on $k$ vectors selects rows $i_1,\dots,i_k$ from the matrix whose columns are
those vectors, and takes the determinant.} The coefficient function is then evaluated at the
point and multiplied on.

Take $\omega := 3\cos(x^1)\,dx^1\wedge dx^3 \in \Omega^2(\mathbb{R}^3)$, the point
$p := (0,1,1)$, and the two vectors
\[v_1 := \begin{pmatrix}1\\0\\0\end{pmatrix}, \qquad
  v_2 := \begin{pmatrix}1\\1\\\tfrac12\end{pmatrix},
  \qquad\text{so}\qquad
  \begin{pmatrix} v_1 & v_2\end{pmatrix} = \begin{pmatrix}
    1 & 1 \\ 0 & 1 \\ 0 & \tfrac12 \end{pmatrix}.\]
The index set is $I = (1,3)$, so we keep rows $1$ and $3$ and discard row $2$:
\[dx^1\wedge dx^3(v_1,v_2) = \det\begin{pmatrix}
  dx^1(v_1) & dx^1(v_2) \\ dx^3(v_1) & dx^3(v_2)\end{pmatrix}
  = \det\begin{pmatrix} 1 & 1 \\ 0 & \tfrac12\end{pmatrix}
  = 1\cdot\tfrac12 - 1\cdot 0 = \tfrac12.\]
The coefficient at $p$ is $3\cos(0) = 3$, and therefore
\[\omega_p(v_1,v_2) = 3 \cdot \tfrac12 = \tfrac32.\]

Two things are worth noticing. The second components of $v_1$ and $v_2$ never entered the
computation, since a form built from $dx^1$ and $dx^3$ is blind to motion in the $x^2$
direction, which is what makes it the natural thing to integrate over a surface in the
$x^1x^3$-plane. And the answer is \emph{signed}: swapping $v_1$ and $v_2$ swaps the determinant's
columns and returns $-\tfrac32$. That sign is exactly the orientation data of
\cref{sec:orientable_manifolds}, and it is why $\int_L\omega$ needs an oriented $L$ to be
well defined.
\end{aiexample}

\begin{ainote}
Applying \cref{def:integral_of_form_over_submanifold} to $\omega_F^{n-1}$ or $\omega_F^2$ from
\cref{ch:differential_forms} recovers exactly the flux integrals of \cref{ch:flux}. Those were,
all along, integrals of forms in this sense. And \cref{def:pullback_of_forms} is what makes the
motivating reparametrization-invariance calculation precise: $\int_L \omega$ does
not depend on the parametrization $\phi$ chosen, only on $L$ and its orientation.
\end{ainote}

% Source: Lukas Krause/Week 12 May 4-May 8 Notes Lukas Krause.pdf, p. 1
% Generator: Gemini 3.6 Flash (High)
\begin{ainote}[Integrating over non-graphical submanifolds via partitions of unity]
\label{note:partitions_of_unity_motivation}
Lukas Krause notes that when a submanifold $\mathcal{M} \subseteq \mathbb{R}^n$ cannot be covered by a single local parametrization chart (for instance, a sphere $S^2$ or a compact manifold with boundary), one cannot directly evaluate $\int_{\mathcal{M}} \omega$ using a single coordinate domain $U$. 

The solution is to decompose any continuous function $f$ on $\mathcal{M}$ into a sum of functions $f = \sum_{l=1}^N \eta_l f$, each supported inside a single coordinate chart domain $U_l$. Integrating each term $\eta_l f$ separately and summing the results yields the total integral over $\mathcal{M}$. The key technical tool for this decomposition is a \newterm{smooth partition of unity} \germanfor{glatte Zerlegung der Eins}.
\end{ainote}

% Source: Jérôme Paschoud/Notizen/Woche 9.2 d-Volumen.pdf, p. 2
% Generator: Gemini 3.1 Pro (High)
\begin{definition}[Support of a function]
\label{def:support_of_function}
Let $U \subseteq \mathbb{R}^n$ be open and $f: U \to \mathbb{R}$ a continuous function. The \newterm{support} \germanfor{Tr\"ager} of $f$ is the closure of the set where $f$ is non-zero:
\[
\operatorname{supp}(f) := \overline{\{x \in U \mid f(x) \neq 0\}}.
\]
If $\operatorname{supp}(f)$ is compact, that is, bounded and with no limit point on the boundary
of $U$, we say that $f$ is \newterm{compactly supported} \germanfor{kompakt getragen}.
% Correction: Claude Opus 5.5 (Extra) --- read "if supp(f) is a bounded set (and supp(f) in U)",
% but the closure is taken in U, so supp(f) lies in U automatically and boundedness alone does not
% give compactness: f = 1 on U = (0,1) has the bounded, non-compact support (0,1).
\end{definition}

% Source: Jérôme Paschoud/Notizen/Woche 9.2 d-Volumen.pdf, p. 2
% Generator: Gemini 3.1 Pro (High)
\begin{example}[Supports of common functions]
\label{ex:supports_functions}
\begin{enumerate}[label=\textbf{(\alph*)}]
  \item For $f: \mathbb{R} \to \mathbb{R}$, $f(x) = \sin(x)$, the function is non-zero on $\mathbb{R} \setminus \pi\mathbb{Z}$. The closure of this set is all of $\mathbb{R}$, so $\operatorname{supp}(f) = \mathbb{R}$. This function is not compactly supported.
  \item Define $g: \mathbb{R} \to \mathbb{R}$ by $g(x) = 0$ for $x \le 0$ and $g(x) = e^{-1/x}$ for $x > 0$. The set where $g(x) \neq 0$ is $(0, \infty)$. Its closure is $[0, \infty)$. Thus $\operatorname{supp}(g) = [0, \infty)$. This is also not compactly supported (since it is not bounded).
\end{enumerate}
\end{example}

% Source: Lukas Krause/Week 12 May 4-May 8 Notes Lukas Krause.pdf, p. 1
% Generator: Gemini 3.6 Flash (High)
\begin{definition}[Standard smooth bump function]
\label{def:smooth_bump_function}
The canonical $C^\infty$ \newterm{bump function} \germanfor{H\"ockerfunktion} $\psi : \mathbb{R}^n \to \mathbb{R}$ supported on the unit ball $B_1(0)$ is defined by
\[
\psi(x) := \begin{cases} \exp\!\left(-\dfrac{1}{1 - |x|^2}\right) & \text{if } |x| < 1, \\[1.5ex] 0 & \text{if } |x| \geq 1. \end{cases}
\]
Since all higher derivatives of $t \mapsto \exp(-1/(1-t))$ approach $0$ faster than any polynomial as $t \to 1^-$, $\psi \in C^\infty(\mathbb{R}^n)$ with $\operatorname{supp}(\psi) = \overline{B_1(0)}$.
\end{definition}

% Source: Lukas Krause/Week 12 May 4-May 8 Notes Lukas Krause.pdf, p. 1
% Generator: Gemini 3.6 Flash (High)
\begin{proposition}[Existence of a smooth partition of unity subordinate to a finite cover]
\label{prop:partition_of_unity_finite}
Let $K \subseteq \mathbb{R}^n$ be compact, and let $\{B_{r_l}(p_l)\}_{l=1}^N$ be a finite open cover of $K$. Then there exist an open set $U \supset K$ and $N+1$ non-negative smooth functions $\eta_1, \dots, \eta_N, \tilde{\eta} \in C^\infty(\mathbb{R}^n, [0,1])$ such that:
\begin{enumerate}[label=\textbf{(\alph*)}]
  \item $\tilde{\eta}(x) + \sum_{l=1}^N \eta_l(x) = 1$ for all $x \in \mathbb{R}^n$;
  \item $\operatorname{supp}(\eta_l) \subseteq B_{r_l}(p_l)$ for each $l \in \{1, \dots, N\}$;
  \item $\operatorname{supp}(\tilde{\eta}) \subseteq \mathbb{R}^n \setminus U$.
\end{enumerate}
In particular, on the neighborhood $U$ of $K$, $\tilde{\eta} \equiv 0$ so that $\sum_{l=1}^N \eta_l(x) = 1$.
\end{proposition}

% Generator: Claude Opus 5.5 (Extra) --- proof added for self-containedness. The idea, "make one
% bump per ball, then divide by their sum", is the whole content; the rest is making sure the
% division never happens where the sum is zero.
\begin{proof}
\textbf{Step 1: shrink the balls.} The function
$x \mapsto \max_l\big(r_l - \lvert x-p_l\rvert\big)$ is continuous and positive on $K$, because
every point of $K$ lies in some $B_{r_l}(p_l)$. On the compact $K$ it therefore has a positive
minimum $\delta$. Every $x \in K$ has some $l$ with $\lvert x-p_l\rvert \leq r_l - \delta$. Discard
the balls with $r_l \leq \delta/2$ and set $\eta_l := 0$ for them: no point of $K$ needs them,
since for those $l$ the inequality would force $\lvert x - p_l\rvert < 0$. For the remaining balls
put $\rho_l := r_l - \delta/2 > 0$ and $\sigma_l := r_l - \delta/4$, so that
$0 < \rho_l < \sigma_l < r_l$ and still $K \subseteq U := \bigcup_l B_{\rho_l}(p_l)$, an open set.
From here on, $l$ ranges over the remaining balls only.

\textbf{Step 2: one bump per ball.} With $\psi$ from \cref{def:smooth_bump_function}, let
$\varphi_l(x) := \psi\big((x-p_l)/\sigma_l\big)$. Each $\varphi_l$ is smooth, non-negative,
positive exactly on $B_{\sigma_l}(p_l)$, and has support $\overline{B_{\sigma_l}(p_l)} \subseteq
B_{r_l}(p_l)$. Their sum $\Phi := \sum_l\varphi_l$ is positive on the compact set
$\overline U = \bigcup_l\overline{B_{\rho_l}(p_l)}$, so $\Phi \geq c$ there for some $c > 0$.

\textbf{Step 3: divide, but only where it is safe.} Choose a smooth $h : \mathbb{R} \to [0,1]$ with
$h(s) = 1$ for $s \leq c/2$ and $h(s) = 0$ for $s \geq c$. For instance, with $g$ from
\cref{ex:supports_functions}\textbf{(b)}, $h(s) := g(c-s)/\big(g(c-s)+g(s-c/2)\big)$; the
denominator never vanishes because at least one of $c-s > 0$ and $s-c/2 > 0$ holds. Now set
\[\tilde\eta := h\circ\Phi, \qquad
  \eta_l := \begin{cases} (1-h(\Phi))\,\dfrac{\varphi_l}{\Phi} & \text{where } \Phi > 0,\\[1ex]
  0 & \text{where } \Phi < c/2.\end{cases}\]
The two cases overlap only where $0 < \Phi < c/2$, and there $1 - h(\Phi) = 0$, so they agree. Each
is smooth on an open set, and the two open sets cover $\mathbb{R}^n$, so $\eta_l$ is smooth. All
values lie in $[0,1]$.
\begin{enumerate}[label=\textbf{(\alph*)}]
  \item Where $\Phi > 0$: $\tilde\eta + \sum_l\eta_l = h(\Phi) + (1-h(\Phi))\cdot\Phi/\Phi = 1$.
    Where $\Phi = 0$: $\tilde\eta = h(0) = 1$ and every $\eta_l = 0$.
  \item $\eta_l$ vanishes wherever $\varphi_l$ does, so
    $\operatorname{supp}\eta_l \subseteq \operatorname{supp}\varphi_l \subseteq B_{r_l}(p_l)$.
  \item On $U$ we have $\Phi \geq c$, hence $\tilde\eta = 0$. So $\{\tilde\eta \neq 0\}$ lies in
    the closed set $\mathbb{R}^n\setminus U$, and so does its closure. \qedhere
\end{enumerate}
\end{proof}

% Generator: Claude Opus 5.5 (Extra) --- the construction in one dimension, without the cut-off h
% (which only matters where both bumps are nearly zero). Two balls B_{1.5}(-0.9), B_{1.5}(0.9)
% cover K = [-1.6,1.6]; bumps phi_l(x) = exp(-1/(1-u^2)), u = (x-p_l)/1.5, peak value e^{-1}
% = 0.368, drawn at vertical scale 4 (peak height 1.47). Bottom panel: eta_l = phi_l/(phi_1+phi_2)
% at vertical scale 1.5, plotted on [-2.3,2.3] where the sum is safely positive; the two curves
% add to 1 (the dashed line at height 1.5) at every x. Below each plot, phi is evaluated with
% max(1-u^2, 0.02) under an indicator so the exponent never overflows outside the ball.
\begin{center}
\begin{tikzpicture}[>=Latex, x=1.5cm, y=1cm,
  declare function={bump(\x,\p)=(abs((\x-\p)/1.5)<1) * exp(-1/max(1-((\x-\p)/1.5)^2, 0.02));
                    etaone(\x)=bump(\x,-0.9)/(bump(\x,-0.9)+bump(\x,0.9));}]
  % top panel: the two bumps
  \draw[->, gray] (-2.7,0) -- (2.8,0) node[right, font=\scriptsize, text=black] {$x$};
  \draw[very thick, OliveGreen] (-1.6,-0.08) -- (1.6,-0.08);
  \node[font=\scriptsize, OliveGreen, below] at (0,-0.1) {$K$};
  \draw[thick, blue!70!black, domain=-2.39:0.59, samples=90, smooth]
    plot (\x, {4*bump(\x,-0.9)});
  \draw[thick, dashed, BurntOrange, domain=-0.59:2.39, samples=90, smooth]
    plot (\x, {4*bump(\x,0.9)});
  \node[font=\scriptsize, blue!70!black] at (-1.75,1.25) {$\varphi_1$};
  \node[font=\scriptsize, BurntOrange] at (1.75,1.25) {$\varphi_2$};
  \node[font=\scriptsize, anchor=west] at (2.45,1.0) {bumps};

  % bottom panel: normalised
  \begin{scope}[yshift=-3.2cm]
    \draw[->, gray] (-2.7,0) -- (2.8,0) node[right, font=\scriptsize, text=black] {$x$};
    \draw[gray, dashed] (-2.3,1.5) -- (2.3,1.5) node[right, font=\scriptsize, text=black] {$1$};
    \draw[very thick, OliveGreen] (-1.6,-0.08) -- (1.6,-0.08);
    \node[font=\scriptsize, OliveGreen, below] at (0,-0.1) {$K$};
    \draw[thick, blue!70!black, domain=-2.3:2.3, samples=90, smooth]
      plot (\x, {1.5*etaone(\x)});
    \draw[thick, dashed, BurntOrange, domain=-2.3:2.3, samples=90, smooth]
      plot (\x, {1.5*(1-etaone(\x))});
    \node[font=\scriptsize, blue!70!black] at (-1.9,1.2) {$\eta_1$};
    \node[font=\scriptsize, BurntOrange] at (1.9,1.2) {$\eta_2$};
    \node[font=\scriptsize, anchor=west, align=left] at (2.0,0.6) {divided by\\ the sum};
  \end{scope}
\end{tikzpicture}
\end{center}


% --- exercises relocated here from the dissolved sheets ---

% Originally: content/week-12.tex, \exercisesheet{12}, problem 12.9
% Source: exercises/Ex12_Analysis2_eng.pdf

\begin{exercise}[Change of variables with differential forms \textnormal{(important)}]
\label{ex:change_of_variables_differential_forms}
Let $U,V \subseteq \mathbb{R}^n$ be open sets, and let $\psi = (\psi^1,\ldots,\psi^n) : U \to
V$ be a diffeomorphism with $\det\!\left(\frac{\partial\psi^j}{\partial x^i}\right) > 0$.
Define an $n$-differential form $\omega$ on $U$ by $\omega = d\psi^1\wedge d\psi^2\wedge\cdots
\wedge d\psi^n$.
\begin{enumerate}[label=\textbf{(\alph*)}]
  \item Show that $\omega = \det\!\left(\frac{\partial\psi^j}{\partial x^i}\right) dx^1\wedge
  \cdots\wedge dx^n$.
  \item From the relation in (a), rewrite the change of variables formula using differential
  forms.
\end{enumerate}
\exinfo{This exercise is Problem 12.9 of Problem Sheet 12, Analysis II, Spring Semester 2026.}
\exsol{sol:change_of_variables_differential_forms}
\end{exercise}

\begin{aiexercise}[Integrating $dx\wedge dy$ over a hemisphere]
% Generator: Claude Opus 5 (High)
\label{ex:ai_area_form_over_hemisphere}
Let $D := \{(u,v) : u^2+v^2 < 1\}$ and let $S^+$ and $S^-$ be the open upper and lower halves of
the unit sphere in $\mathbb{R}^3$, each oriented by the \emph{outward} normal. Consider the
$2$-form $dx\wedge dy$ on $\mathbb{R}^3$.
\begin{enumerate}[label=\textbf{(\alph*)}]
  \item Parametrize $S^+$ as a graph, $\phi(u,v) := \big(u,\ v,\ \sqrt{1-u^2-v^2}\big)$ over
    $D$. Compute $\phi^*(dx\wedge dy)$ and deduce $\int_{S^+} dx\wedge dy$.
  \item Do the same for $S^-$ with $\psi(u,v) := \big(u,\ v,\ -\sqrt{1-u^2-v^2}\big)$. The
    pullback is the same as in \textbf{(a)}, yet the answer is not. Which hypothesis of
    \cref{def:integral_of_form_over_submanifold} accounts for the difference?
  \item Deduce $\int_{S^2} dx\wedge dy$, and confirm the value a second time by writing
    $dx\wedge dy = d(x\,dy)$ and applying \cref{thm:stokes_general}.
\end{enumerate}
\exsol{sol:ai_area_form_over_hemisphere}
\end{aiexercise}
